%----------------------------------------------------------------------------------------
%	PACKAGES AND OTHER DOCUMENT CONFIGURATIONS
%----------------------------------------------------------------------------------------
\documentclass[a4paper,12pt]{article}
\usepackage[english]{babel}
\usepackage[latin1]{inputenc}
\usepackage{amsmath}
\usepackage{xcolor}
\usepackage{amssymb}
\usepackage{amsfonts}
\usepackage{graphicx}
\usepackage{geometry}
\definecolor{link}{rgb}{0.4, 0.2, 0.2}
\geometry{
	a4paper,
	total={170mm,257mm},
	left=20mm,
	top=20mm,
}
\usepackage{hyperref}
\hypersetup{
	unicode=false,          % non-Latin characters in Acrobat’s bookmarks
	pdftoolbar=true,        % show Acrobat’s toolbar?
	pdfmenubar=true,        % show Acrobat’s menu?
	pdffitwindow=false,     % window fit to page when opened
	pdfstartview={FitH},    % fits the width of the page to the window
	pdftitle={Research Planning Report},    % title
	pdfauthor={Aaron Panaitescu},     % author
	pdfkeywords={neuroplasticity, Parkinson's, stimulation}, % list of keywords
	pdfnewwindow=true,      % links in new PDF window
	colorlinks=true,       % false: boxed links; true: colored links
	linkcolor=link,          % color of internal links (change box color with linkbordercolor)
	citecolor=link,        % color of links to bibliography
	filecolor=link,      % color of file links
	urlcolor=link           % color of external links
}



% \makeglossaries
% \newcommand{\ac}[1]{\pdftooltip{\gls*{#1}}{\glsentrylong{#1}}}
% \newacronym{abc}{ABC}{A B and C}

\begin{document}
\begin{titlepage}
	\newcommand{\HRule}{\rule{\linewidth}{0.5mm}}
	\setlength{\topmargin}{0in}
	\center

	\begin{flushleft} \large
		\begin{minipage}{0.4\textwidth}
			\includegraphics[scale=0.14]{imperial.png}
		\end{minipage}
	\end{flushleft}
	\vspace{4cm}

	%----------------------------------------------------------------------------------------
	%	TITLE SECTION
	%----------------------------------------------------------------------------------------
	\textbf{\large Literature review and thesis proposal}\\[0.1cm]
	{\large MRes. Neurotechonlogy}\\[0.5cm]

	\HRule \\[0.4cm]
	{\Large \bfseries Investigating plasticity in \\ Cortico-Basal Ganglia-Thalamus models \\ to improve stimulation-based treatments }
	\HRule \\[1cm]

	%----------------------------------------------------------------------------------------
	%	AUTHOR SECTION
	%----------------------------------------------------------------------------------------

	{\large Aaron Panaitescu \\
	(\textit{CID:} 02054726) \\[0.4cm]
	\textbf{Supervisor:} \\
	Hayriye Cagnan}

	%----------------------------------------------------------------------------------------
	%	DATE SECTION
	%----------------------------------------------------------------------------------------

	\vfill
	{Word Count: ????}\\[0.4cm]
	{\large \today}\\[0.8cm]
\end{titlepage}
%----------------------------------------------------------------------------------------
%	REPORT CONTENT
%----------------------------------------------------------------------------------------
\tableofcontents
\newpage

\section{Introduction / Abstract?}
**Stylistically, I would prefere to do an abstract here and breakdown the concepts in the following section**
\section{Background and literature review}
**start this off nicely with a diagram of the research gap**

\subsection{Parkinson's Disease}
\textit{Outline} \cite{del2018advances}

** 1. Loss of SNc dopaminergic neurons.**

** 2. indirect GPe $\rightarrow$ STN pathway $\uparrow$,
hyperdirect Cortex $\rightarrow$ STN pathway $\downarrow$.
(dimmer switch model \cite{helmich2012cerebral}, \cite{west2022stimulating})**

** \textbf{add diagram} **

** 3. Hypersynchrony in the Basal Ganglia.**

\subsection{DBS: theory and practice}
**DBS as the state of the art in treatment.**

**Limitations of DBS (invasiveness, side effects, it needs to be on permenently,
why 130Hz? when tremors are $\sim$ 20Hz)**

$\Rightarrow$ plenty of things to be improved

**citations needed**

**also should probably add a bit on other types of stimulation**

(e.g. tACS \cite{saturnino2017target} \cite{schwab2020spike} \cite{schwab2019modulation})


\subsection{Stimulating at the right time?}
**important** \cite{cagnan2017stimulating} \cite{beudel2018adaptive} \cite{west2022stimulating}

\subsection{Plasticity to recover network states}
**mention** \cite{lebedev2017brain} \cite{cramer2011harnessing}

\subsection{Neuron-level vs. Mean-field models}
**briefly, in general and expand in the context of plasticity**

**cover** \cite{jansen1995electroencephalogram} (\cite{hodgkin1952measurement} does this really need to be cited?)

**important** \cite{terman2002activity} \cite{rubin2012basal} \cite{duchet2023mean} \\
\cite{shupe2021integrate} \cite{schwab2020spike}

\subsection{**Other ways of improving stimulation-based treatments**}
**stimulation parameter optimizations, closing the loop (e.g. aDBS \cite{beudel2018adaptive})**
\textit{
	Maybe this can be folded into the stimulating at the right time part
	since they are pretty closely related. Also should mention the idea that
	these methods are not mutually exclusive, meaning that in principle they
	could be combined, adaptive stim closing the loop, with inducing plastic
	changes as target
}

\section{Project Plan}


\subsection{Aims}
\begin{enumerate}
	\item Model neuroplasticity in a Parkinsonian CGBT network
	\item Investigate the viabillity of harnessing plasticity to remove the
	      system from the pathological state and analyze the dynamics that follow

	      **here i care about things like for how long and how does the network change.
	      To what degree can we induce changes etc.**

	      **should look into viable timescales + noise analyze how long does it take to rebound to
	      pathological state**
	\item Try to link potential results to potential stimulation protocols?
\end{enumerate}

\subsection{Methodolgy} **How indeed?** HH/IF Pakrkinsoni model + plasticity rules, trying
different stimulation-based protocols (link with experimental data?)

**look into how to tune IF model parameters from experimental data**

\subsection{Timeline}
**Gant chart thingy**



%**********************************************%
\newpage
%----------------------------------------------------------------------------------------
%	REFERENCES
%----------------------------------------------------------------------------------------
\addcontentsline{toc}{section}{References}
\bibliography{refs}
\bibliographystyle{apalike}
\newpage

%----------------------------------------------------------------------------------------
%	APPENDIX  
%----------------------------------------------------------------------------------------
\appendix
\section*{Appendices}
\addcontentsline{toc}{section}{Appendices}
\renewcommand{\thesubsection}{\Alph{subsection}}
\subsection{First appendix}
\subsection{Second appendix}

\end{document}
